\section{Traditional Optical Character Recognition Methods}

Traditional Optical Character Recognition (OCR) methods, developed primarily before the advent of modern machine learning and deep learning techniques, rely heavily on rule-based, template-based, or heuristic-driven approaches to transform text from images into machine-readable formats. These methods leverage image processing techniques, handcrafted feature extraction, and predefined geometric patterns to recognize characters, often with a focus on digitized scanned documents. While foundational to the field of automated text recognition, these approaches are characterized by their dependence on manual effort and domain-specific expertise, which limits their scalability and adaptability in complex visual environments. This section reviews the key paradigms in traditional OCR, traces their evolution over time, and highlights their contributions and shortcomings.

\subsection{Template Matching and Early Approaches}

One of the earliest notable efforts in traditional OCR is the work of Mori, Suen, and Yamamoto \cite{Mori1992}, who synthesized and developed first-generation OCR techniques. The approach during this period utilized basic image processing to perform template matching, comparing images pixel-by-pixel against a predefined library of character templates. While straightforward and computationally efficient for uniform, typewritten documents, this method was highly sensitive to variations in font style, size, and image quality, making the system brittle when confronted with diverse document conditions.

Lenet and colleagues later pioneered the application of convolutional neural network precursors combined with traditional sliding-window classifiers for digit recognition \cite{LeCun1998}. This hybrid approach improved upon purely template-based systems by incorporating learned local features, but still required carefully engineered preprocessing pipelines and struggled to generalize beyond constrained digit vocabularies. Casey and Lecolinet \cite{Casey1996} further systematized recognition strategies through a taxonomy of structural decomposition approaches, in which characters are broken into primitive strokes and reassembled via graph-matching algorithms. This approach excelled in domain-specific contexts such as postal code recognition, but its reliance on a fixed set of structural primitives restricted flexibility when confronted with diverse font families, as adapting the system to new scripts required redesigning the primitive library from scratch.

\subsection{Structural and Feature-Based Methods}

The OCR field advanced considerably with research on structural and topological feature extraction, exemplified by the work of Tappert et al. \cite{Tappert1990}. Rather than relying on pixel-level comparison, this approach extracted key geometric characteristics from text, such as loops, lines, and stroke intersections, to serve as inputs to classification algorithms. This strategy proved particularly valuable for processing standardized printed documents, where character geometry is well-defined and consistent. However, it required significant manual tuning to design features robust to noise and font variation, highlighting the heavy dependence of traditional methods on human intervention and domain knowledge.

Nagy \cite{Nagy2000} extended this paradigm by exploring document layout analysis as a precursor to character recognition. The proposed system uses geometric segmentation to identify text blocks, columns, and lines before passing individual regions to a character classifier, marking a step forward in handling multi-column documents. Its dependence on clean binarization and rigid layout assumptions, however, limited robustness across varied real-world document conditions, where uneven illumination, physical page curvature, and background clutter frequently violate the assumptions required for reliable segmentation.

\subsection{Statistical and Sequence-Based Methods}

A significant conceptual advance came with the application of Hidden Markov Models (HMMs) to text recognition, as demonstrated in works such as those by Anigbogu \cite{Anigbogu1995}. Rather than treating each character as an isolated unit, HMM-based systems model text as a sequence of observations derived from sliding vertical windows across a text line, allowing recognition to proceed without explicit character segmentation. This formulation was groundbreaking because it enabled the recognition of cursive scripts and noisy text in which individual character boundaries are ambiguous. The effectiveness of HMM-based systems, however, depended heavily on the quality of manually crafted sliding-window features and the accompanying statistical language model, limiting their transferability across scripts and document types.

\subsection{Traditional Vietnamese OCR}

Researchers further expanded the scope of OCR by addressing highly complex writing systems, with Vietnamese receiving specific attention due to its dense diacritical structure \cite{Le2025}. Traditional approaches to Vietnamese OCR typically combined rule-based projection profiles to separate main characters from tonal marks, followed by heuristic scoring methods to reassociate diacritics with their base characters. This pipeline exploited the relatively predictable vertical arrangement of Vietnamese composite glyphs in printed documents, where tone marks appear consistently above or below the base vowel region.

While functional for clean, high-resolution administrative documents, these methods struggled to scale across varied background noise conditions and failed almost entirely on scene text images captured in natural environments. The projection-profile assumption breaks down when text lines are skewed, when characters overlap due to tight tracking, or when tonal marks are partially obscured, all of which are common in real-world Vietnamese document collections. The fragility of this approach under these conditions directly motivated the adoption of data-driven recognition methods capable of learning robust representations of Vietnamese composite glyphs without hand-engineered segmentation rules.

\subsection{Strengths and Limitations}

Traditional OCR methods possess notable strengths in controlled settings. Their rule-based and template-driven design ensures predictable, interpretable outputs well-suited to digitizing typewritten administrative records or standardized printed forms. They perform reliably on uniform, high-resolution documents with consistent fonts and layouts, and require minimal labeled data and computational resources compared to modern deep learning systems, making them appropriate for early digitization efforts or resource-limited environments.

Their limitations, however, are fundamental rather than incidental. Handcrafted features and rigid templates severely restrict scalability and adaptability across diverse document types. HMM-based systems are brittle under varying noise conditions due to their dependence on manually crafted features. Perhaps most critically, traditional methods lack the capacity to model contextual dependencies across characters or words, producing recognition outputs that are locally plausible but globally inconsistent, a particularly serious limitation for morphologically complex scripts such as Vietnamese in which diacritical context determines word identity. These shortcomings collectively drove the transition toward automated, data-driven approaches, which are examined in the following section.
